\qznuappendix{附录 B　“世遗泉州”项目学习任务单}

\section*{一、小组任务清单}

\begin{table}[htbp]
  \centering
  \caption{“世遗泉州”项目小组任务清单}
  \label{tab:tasklist}
  \begin{tabularx}{\textwidth}{@{}L{2.6cm}Yp{2.2cm}@{}}
    \toprule
    任务 & 具体要求 & 完成时间 \\
    \midrule
    选定遗产点 & 从泉州世界遗产名录中选择一处，说明选择理由 & 第1周 \\
    资料卡片 & 至少收集三类资料，标注来源并判断可靠性 & 第2周 \\
    现场观察 & 完成观察记录表，记录建筑细节、人物活动与环境信息 & 第3周 \\
    人物访谈 & 访谈至少一位讲解员、居民或游客，整理访谈要点 & 第3周 \\
    解说词初稿 & 面向指定受众，完成800字左右的解说词 & 第4周 \\
    同伴互评 & 依据评价量表提出两条具体修改建议 & 第5周 \\
    最终作品 & 完成手册、音频、短视频或展览作品 & 第7周 \\
    展示反思 & 参与公开展示，提交反思日志 & 第8周 \\
    \bottomrule
  \end{tabularx}
\end{table}

\section*{二、资料筛选记录表}

\begin{table}[htbp]
  \centering
  \caption{资料筛选记录表示例}
  \label{tab:source}
  \begin{tabularx}{\textwidth}{@{}L{2.2cm}L{3.2cm}Yp{2.2cm}@{}}
    \toprule
    资料来源 & 主要信息 & 可靠性判断 & 拟用位置 \\
    \midrule
    官方网站 & 遗产点基本信息、保护范围 & 权威性较高，可作为事实依据 & 解说词开头 \\
    地方文献 & 历史沿革、人物故事 & 需核对出版信息与作者背景 & 主体部分 \\
    现场访谈 & 讲解员经验、居民记忆 & 属于口述材料，需注明访谈对象与时间 & 细节描写 \\
    网络文章 & 旅游攻略、传说故事 & 可靠性不一，需与权威资料交叉验证 & 导入或延伸 \\
    \bottomrule
  \end{tabularx}
\end{table}

\section*{三、反思日志提纲}

请围绕以下问题撰写不少于300字的反思日志：你在项目中承担了什么任务？你遇到的最大困难是什么？你是如何解决的？你对泉州世界遗产的认识发生了什么变化？如果再来一次，你会在哪些方面改进？你认为一篇好的解说词应该具备哪些特点？
